%%
%% 2020 - 08 -03 φ
%%

\subsubsection{Erwartungswert}\index{Erwartungswert}


\begin{definition}{Zufallsvariable}{}
	Eine \textbf{Zufallsvariable}\index{Zufallsvariable} $X(\omega)$ ordnet jedem Ergebnis eine Zahl
	(typischerweise den Gewinn oder Verlust) zu.
\end{definition}


\begin{definition}{Erwartungswert}{}
	Mit $\mu$ wird der \textbf{Erwartungswert}\index{Erwartungswert} bezeichnet.
	$\mu$ ist das gewichtete arithmetische Mittel, das die Zufallsvariable
	annimmt.
\end{definition}



\begin{beispiel}{Erwartungswert}{}
  Bei einem Spiel sind in einer Urne drei rote und eine grüne
  Kugel. Der Spieleinsatz sind CHF 1.-. Nun darf eine Kugel gezogen
  werden. Es gelten die folgenden Regeln:
  \begin{itemize}
  \item «rot» gilt als Verlust und der Einsatz (CHF 1.-) ist verloren
  \item «grün» gilt als Gewinn und man erhält CHF 6.- minus den Einsatz
  \end{itemize}
  Die Zufallsvariable sind also wie folgt zu definieren:
  \begin{itemize}
  \item $X(\text{«rot»}) = \LoesungsRaum{-1}$
  \item $X(\text{«grün»}) = \LoesungsRaum{+5}$
  \end{itemize}

  Weil wir aber viel häufiger «rot» ziehen werden, erwarten wir bei
  vier Spielen drei Verluste und einen Gewinn. Bei vier Spielen ist
  (sofern jede Kugel genau einmal gezogen wurde) unser Gewinn

  $$ \mu = \LoesungsRaumLen{50mm}{( -1 -1 -1 + 5) / 4} = \LoesungsRaumLen{50mm}{\frac{3\cdot{}(-1) + 1\cdot{}5}{3+1}}
  = ... = \LoesungsRaumLen{12mm}{0.50} \text{ CHF}$$
\end{beispiel}

Allgemein wird der Erwartungswert $\mu$ wie folgt berechnet:

\begin{gesetz}{Erwartungswert}{}\index{Erwartungswert!Formel}
$$\mu = \sum P(\omega_i)\cdot{} X(\omega_i)$$
\end{gesetz}
\newpage


\subsection*{Referenzaufgabe (optional)}

Ich schlage Ihnen folgendes Würfelspiel vor.

Sie Würfeln und erhalten dabei die Quadratzahl der erzielten
Würfelpunkte in CHF ausbezahlt. Wenn Sie \zB eine \epsdice{5} werfen,
so erhalten Sie CHF 25.-.

Natürlich müssen Sie einen Grundeinsatz bezahlen, damit Sie bei meinem
Spiel überhaupt mitmachen dürfen.

Drehen wir nun die Rollen einmal um und betrachten Sie sich selbst als
Betreiber dieses «Casinos», also dieses Spiels.

Wie viel CHF soll der Spieler pro Wurf als Einsatz bezahlen, sodass
dieses Spiel mit minimaler Wahrscheinlichkeit besser für den
Casino-Betreiber, also für Sie ausgeht?

\TNTeop{
  Dazu berechnen wir zuerst das, was wir erwarten, wenn dieses Spiel
  oft gespielt wird. Jedes Würfelresultat wird mit $\frac16$
  Wahrscheinlichkeit auftreten. Somit muss jeder Gewinn ja nur mit
  $\frac16$ Wahrscheinlichkeit ausbezahlt werden.

  Wenn wir nun alle Rückzahlungen mit $\frac16$ multiplizieren und
  diese CHF-Werte zusammenzählen, erhalten wir

  $$\frac16\cdot{} (1 + 4 + 9 + 16 + 25 + 36) = 15.166... CHF$$

  Diese Zahl nennen wir den \textbf{Erwartungswert}.

  In obigem Beispiele ist $\mu = 15.166...$ CHF und die Zufallsvariable
  ordnet jedem Ergebnis (1..6) seinen auszuzahlenden Wert zu.

  Zufallsvariable hier:

  $$X(\omega) = 1.- \text{ für } \epsdice{1}; 4.- \text{ für }
  \epsdice{2} ; ...$$
  
  Wenn wir als Einsatz CHF 16.- verlangen, wird sich dieses Spiel im
  Endeffekt positiv für das «Casino» auswirken. Wenn wir hingegen nur
  CHF 15.- als Einsatz verlangen, wird früher oder später der Spieler
  mehr auf dem Konto haben.
}%% END TNTeop
%%\newpage


\subsection*{Erwartungswert bei Binomialverteilungen}

\begin{gesetz}{Erwartungswert bei Binomialverteilungen}{}
  Bei Binomialverteilungen mit Wahrscheinlichkeit $p$ und
  Stichprobenumfang $n$ kann der Erwartungswert $\mathbb{E}(X)$ wie folgt
  berechnet werden:
  $$\mathbb{E}(X) = n\cdot{}p$$
  Anstelle von $\mu$ schreiben wir bei Binomialverteilungen in $X$ auch $\mathbb{E}(X)$.
\end{gesetz}


\subsection*{Referenzaufgaben}


Angenommen, die Aussage des Pharma-Herstellers sei korrekt, so gehen
wir stets von der \textbf{Nullhypothese} aus:

«Das Medikament verspricht bei 80\,\% der Patienten Linderung.»

\aufgabenFarbe{Wie groß ist der Erwartungswert $\mathbb{E}(X)$ in einer Stichprobe
von 100 Probanden, die das Medikament mit 80 \% Linderungsversprechen
zu sich nehmen, dass tatsächlich Linderung eintrifft?}

\TNT{2}{$$\mathbb{E}(X) = n\cdot{}{} = 100\cdot{}0.8 = 80 \text{ Personen}$$}

\aufgabenFarbe{Basil Ballisti der Basketballspieler trifft mit einer
Wahrscheinlichkeit von 85 \%. Er wirft 20 Bälle. Was ist sein
Erwartungswert $\mathbb{E}(X)$, wenn $X$ die Anzahl der Treffer sein
soll?}

$$\mathbb{E}(X) = \LoesungsRaumLen{80mm}{n\cdot{}p = 20 \cdot{} 0.85 = 17 \text{ Treffer}}$$

Basil Ballisti wird mit der höchsten Wahrscheinlichkeit um die
\LoesungsRaumLen{30mm}{17} Treffer landen.
\newpage
